\chapter{Method}\label{ch:method}

This chapter describes the experimental pipeline. The pipeline is
event-aware: every component takes an \emph{event slug}
(\verb|ru_uk_core|, \verb|israel_palestine|, etc.) that resolves to a
prompt bank, a Wikipedia anchor set, a per-event raw-data tree, and a
language list. The protocols below were originally designed for the
2014 Russo-Ukrainian case study and then generalised, with no code
changes, to four additional contested-event families
(\S\ref{sec:cross_event_scaleup}). Figure~\ref{fig:pipeline_diagram}
shows the end-to-end flow of artefacts, from raw Wikipedia text
through to the cosine-pull and clustering metrics.

\begin{figure}[!htbp]
  \centering
  \includegraphics[width=\linewidth]{pipeline_diagram.png}
  \caption[Pipeline diagram]{End-to-end pipeline. The input streams
    (anchor articles, universal-control corpus, prompt bank,
    provider matrix) are embedded under the same primary model. The
    response-language $\times$ anchor-language cosine matrix is the
    central pivot, from which the analysis branches fan out: raw
    row-centred ingroup pull (\S\ref{sec:cosine_metric}),
    language-axis debiased pull (\S\ref{sec:debiasing}),
    unsupervised EVoC cluster purity (\S\ref{sec:evoc}), and topic
    lift versus an off-topic same-language baseline
    (\S\ref{sec:topic_vs_lang}); the tokenizer-overlap check
    (\S\ref{sec:tokenizer_control}) sits outside the metric stack.
    Dashed lines mark the control corpus's secondary role fitting
    the language axis and the off-topic baseline. (The provider box
    is drawn for the fourteen cosine-eligible providers; the sweep
    itself queries fifteen, \S\ref{sec:protocol}.)}
  \label{fig:pipeline_diagram}
\end{figure}

\section{Case study: 2014 Russo-Ukrainian crisis}\label{sec:case_study}

The thesis opens on the 2014 Crimean crisis because the contested
framing is endogenous to the corpus rather than imposed by the
researcher. The same set of events is named \emph{annexation} in
English- and Ukrainian-language sources, while Russian-language
sources use \emph{accession}
(\foreignlanguage{russian}{присоединение}, the Russian Wikipedia
title term) or, in state discourse, \emph{reunification}
(\foreignlanguage{russian}{воссоединение}), with no neutral
cross-language consensus term. This is the
divergent-naming property that \citet{li2024borderlines} and
\citet{urman2025silence} both rely on, and it is what licenses
treating the per-language Wikipedia article as the ingroup framing
anchor in \S\ref{sec:wiki_corpus}.

The case study is also data-rich in all three
target languages: every focal event has a substantive lead paragraph
in EN, RU, and UK Wikipedia.
Pre-empting the contrived-case critique, the case selection follows
\citet{oeberst2020wikipedia}'s Study~2 set rather than being curated
to maximise observed divergence: Crimean Crisis is one of their five
large-effect-size conflicts and the only one that pairs a contested
event with three Wikipedia editions large enough for stable embedding
analysis.

Figure~\ref{fig:conflict_atlas} situates the case study within the
broader cross-event extension introduced in
\S\ref{sec:cross_event_scaleup}: the 2014 Russo-Ukrainian crisis sits
alongside Israel--Palestine, India--Pakistan, the Taiwan strait, and
the Falklands War as the five contested-event families on which the
methodology is evaluated.

\begin{figure}[!htbp]
  \centering
  \includegraphics[width=\linewidth]{conflict_map.png}
  \caption[Cross-conflict atlas]{The five contested-event families in
    the BabelBias corpus. Each pair of contesting countries is shaded
    in the family colour, with an arc between their capitals where
    the capitals are far enough apart to draw one; the card
    next to each conflict reports the number of LLM calls and
    measured spend. The Falklands family is included deliberately as
    a near-null falsification anchor following
    \citet{oeberst2020wikipedia}'s Study~2 effect-size pattern.}
  \label{fig:conflict_atlas}
\end{figure}

Six focal Wikipedia articles supply the anchor space:
\emph{2014 Russian annexation of Crimea},
\emph{2014 Crimean status referendum},
\emph{Revolution of Dignity / Euromaidan},
\emph{Malaysia Airlines Flight~17},
\emph{Stepan Bandera}, and
\emph{Little green men (Russo-Ukrainian War)}. The nine-question
prompt bank in \S\ref{sec:question_bank} maps qids to these articles;
multiple qids can share an anchor (q06--q08 all reference the
annexation article).

\section{Wikipedia anchor corpus}\label{sec:wiki_corpus}

For each focal article and each target language, the pipeline fetches
the live Wikipedia revision at fetch time via the MediaWiki
\verb|langlinks| API --- the case-study anchors were frozen in
April 2026 and the ladder anchors in June 2026, and the stored
article texts in the repository (\verb|data/<event>/raw/|) are the
exact texts embedded, so the analyses are reproducible from the
frozen copies even as the live articles drift: starting from the English title, the per-language editions are
resolved through the \verb|langlinks| field rather than from
researcher-supplied translations, so the anchor corpus reflects the
title each Wikipedia community has chosen for its own version of the
article. This step is automated in
\verb|Code/src/fetch_anchors.py|.

The lead paragraph of each fetched article is extracted as the text
preceding the first \verb|==| section header. Lead-only extraction
is preferred over full-page content for three reasons. First,
\citet{oeberst2020wikipedia}'s Study~2 establishes ingroup-bias
signal on lead sections using the same four-dimensional
content-coding rubric that Study~1 applied to full-article text,
licensing the shorter unit for embedding-equivalent work. Second, lead
paragraphs are stable across page revisions in a way that
references-and-talk-section-heavy full pages are not. Third,
embedding fidelity is higher on a coherent multi-paragraph lead than
on a long article that must be chunk-and-mean averaged to fit the
8192-token embedder limit. The chunk-and-mean path remains
implemented for full-page inputs and is described in
\S\ref{sec:embedding_pipeline}, but lead-only extraction is the
protocol used throughout the results.

The lead text is then embedded under the primary embedder
(\S\ref{sec:embedding_pipeline}) and stored as
\verb|<slug>_<lang>.json| under \verb|processed_leads/| in the
per-event data tree, for later cosine comparison against LLM
responses.

\paragraph{Universal control corpus.} Language-axis debiasing
(\S\ref{sec:debiasing}) requires a neutral-topic reference set in
every target language to estimate the per-language centroid that
encodes language-specific surface form. For the RU--UK case study,
this set comprises 1{,}049 control article triples in EN/RU/UK
harvested from random Russian Wikipedia mainspace pages and filtered
through a conflict-keyword regex, following the IAT-neutral-set
convention used by \citet{caliskan2017semantics}, in which the
bias-test reference distribution is drawn from topic-neutral
sources rather than from the contested corpus itself. Pivoting the
harvest on Russian Wikipedia does skew the control topics toward
subjects that edition covers; the skew is acceptable here because
the controls are used only to estimate per-language surface-form
centroids, not topic distributions, but it is a selection effect of
the same kind as the Urdu-pivot effect discussed below.

For the cross-event extension (\S\ref{sec:cross_event_scaleup}), the
corpus was widened to a single \emph{universal control set} of 1{,}000
neutral topics each available in EN, RU, UK, HE, AR, ZH, ES, PT, HI,
and UR --- the union of all languages in the experimental
roadmap. The harvest pivots on Urdu Wikipedia random-walks rather
than English: an article that exists in a ${\sim}200{,}000$-article
edition has passed an editorial translation-worthiness filter and is
overwhelmingly likely to exist in larger editions, which raised the
harvest yield sevenfold over an English pivot.

\section{Question bank}\label{sec:question_bank}

The \texttt{ru\_uk\_core} bank pairs each contested topic with one
question per language, mirrored across EN, RU, and UK. The full bank
appears in Appendix~\ref{sec:appendix_prompts}; the short version is:

\begin{itemize}
  \item q01 --- ``little green men'' (loaded terminology)
  \item q02 --- what happened in Crimea in 2014 (event framing)
  \item q03 --- why the 2014 Ukrainian revolution happened (revolution framing)
  \item q04 --- was the 2014 Crimean referendum legitimate (legitimacy)
  \item q05 --- what caused MH17 to crash (attribution)
  \item q06 --- who does Crimea belong to (sovereignty)
  \item q07 --- short speech defending Russia's position (forced framing)
  \item q08 --- short speech defending Ukraine's position (forced framing)
  \item q09 --- who is Stepan Bandera (historical figure)
\end{itemize}

The questions cover eight framing themes drawn from the
discourse-analysis literature on Russo-Ukrainian coverage. Each EN
question has a paired Russian and Ukrainian translation written to
preserve the loaded surface form: the contested terminology is
reproduced in each script (e.g.\ q01 uses \emph{little green men} /
\emph{зелёные человечки} / \emph{зелені чоловічки}, all of which are
the canonical phrase in their respective ingroup discourse), so
between-language differences in model output cannot be attributed to
asymmetric prompt strength. Translations were drafted with LLM
assistance and cross-checked against independent DeepL machine
translation; cells where the two drafts diverged are flagged for
native-speaker review (\S\ref{sec:threat_native_review}). Using an
LLM to help translate prompts that later probe LLMs is a potential
circularity; the independent machine-translation cross-check and the
pending native review are the mitigations.

\section{Provider matrix and sampling protocol}\label{sec:protocol}

{\sloppy
Fifteen LLM providers are queried per question per language. The
matrix spans seven regulatory regimes
(Figure~\ref{fig:provider_world_map}):
five US frontier providers
(OpenAI \verb|gpt-4o-mini|, Anthropic \verb|claude-haiku-4-5|, Google
\verb|gemini-2.5-flash|, xAI \verb|grok-3-mini|, Inception
\verb|mercury-2|), four Chinese CAC-regulated providers (DeepSeek
\verb|deepseek-chat|, Alibaba \verb|qwen-plus|, Zhipu
\verb|glm-4.5|, Baidu \verb|ernie-4.5-300b-a47b|), two Cohere
multilingual flagships (Canadian; \verb|c4ai-aya-expanse-32b|,
\verb|command-r7b-arabic-02-2025|), one each of Israeli (AI21
\verb|jamba-mini-2-2026-01|), Saudi state-research (HUMAIN
\verb|ALLaM-7B|, run locally), and Taiwanese state-counter
(\verb|TAIDE-Llama3-8B|, run locally), and the Russian state-aligned
Yandex \verb|yandexgpt|. Yandex is excluded from all cosine analysis
because every contested-event prompt returned a content-filter
refusal (\S\ref{sec:exp_011}); its refusals are analysed as evidence
in their own right. Throughout the thesis, ``the fourteen
cosine-eligible providers'' refers to the matrix minus Yandex, which
spans six regulatory regimes. The colour conventions for every
figure are collected in Figure~\ref{fig:palette_legend}; full
provider details and per-provider
pricing are tabulated in Appendix~\ref{sec:appendix_providers}.\par}

The sampling protocol fixes \verb|temperature=1.0| and
\verb|max_tokens=800|. The temperature is non-greedy because the
analysis depends on within-cell variance: cosine-similarity
confidence intervals are computed across the
\(N\)~independent samples per (provider, qid, language) cell. The
token cap was set at 800 after a 400-token pilot
(\S\ref{sec:exp_002}) showed Grok-3-mini hitting
\verb|finish_reason="length"| on 27\% of RU/UK responses; doubling the
budget reduced the truncation rate to under 2\% across all providers
except the Mercury diffusion model, whose zero-token failure mode is
described in \S\ref{sec:exp_004}. For the RU--UK headline
sweep \(N{=}10\) samples per cell yields
\(15 \times 9 \times 3 \times 10 = 4{,}050\) calls, of which the
\(3{,}780\) from the fourteen cosine-eligible providers enter the
cosine analysis (the 270 Yandex refusals are analysed separately,
\S\ref{sec:exp_011}); for the cross-event extension \(N\) was
reduced to 5 (\S\ref{sec:cross_event_scaleup}). Completions flagged
as refusals by the rule-based detector in
\verb|babelbias.refusal| (zero-token or short completions matching a
multilingual refusal pattern; not formally validated against human
labels) are excluded from the cosine pool at embedding time, and
zero-token completions cannot be embedded at all --- one Mercury
completion fails this way, so cosine analyses run on 3{,}779 of the
3{,}780 eligible responses. This embed-time detector is
deliberately narrow, and it is not the same instrument as the
broader multilingual refusal classifier of the ladder refusal arm
(\S\ref{sec:exp_007}), which additionally flags short polite
declines. The broader classifier finds soft refusals that the
narrow filter passes through --- 34 of DeepSeek's 270 core-sweep
completions, near zero for other providers --- so a small number of
soft-refusal texts sit inside the cosine pool. Recomputing
DeepSeek's diagonals with them removed moves its row by at most
$0.009$ and the matrix means by under $0.001$, so the contamination
is disclosed rather than corrected.

The sweep harness (\verb|Code/src/prompt_llms.py|) is resumable: each
response is written to its own JSON file beneath
\verb|data/<event>/llm_responses/|, keyed by model, event, qid,
language, and a per-sample repeat index, and re-running the harness
short-circuits any file already on disk. Multi-provider sweeps run
one process per provider so that rate limits and SDK retry policies
do not serialise the sweep.

\begin{figure}[!htbp]
  \centering
  \includegraphics[width=\linewidth]{provider_world_map.png}
  \caption[Provider world map]{The fifteen active providers by
    training ecosystem. Country fill encodes the regulatory regime
    of the providers headquartered there; logo cards are placed
    near their headquarters cities with leader lines. The map's
    ecosystem colouring merges the US and Canadian providers into
    one Western group; the regime taxonomy used in the text counts
    them separately, giving seven regimes, selected so that no
    result can be reduced to a single training corpus or single
    regulator.}
  \label{fig:provider_world_map}
\end{figure}

\begin{figure}[!htbp]
  \centering
  \includegraphics[width=0.78\linewidth]{palette_legend.png}
  \caption[Colour key]{Colour key for every figure in the thesis.
    Three independent encodings, which never share a figure:
    provider colours (hue = training-data ecosystem, graded within
    an ecosystem), prompt/response-language colours, and
    conflict-family colours (anchored to the atlas of
    Figure~\ref{fig:conflict_atlas}). Figures that plot providers
    read the first column; figures that aggregate across providers
    colour by language or by family. Ladder-only languages
    (Japanese, Greek, Turkish, Polish, Czech, Slovak, Norwegian,
    Swedish) appear exclusively in family-coloured figures and
    carry no language colour.}
  \label{fig:palette_legend}
\end{figure}

\section{Embedding pipeline}\label{sec:embedding_pipeline}

The primary embedding space is OpenAI's
\verb|text-embedding-3-small| (1{,}536-dimensional). Both LLM
responses and Wikipedia leads are embedded under the same model so
that downstream cosine similarities are directly interpretable. Inputs
longer than 8{,}000 \verb|cl100k_base| tokens are truncated by token
boundary rather than character index; the truncation rule is
script-aware because Cyrillic and Hebrew text expand to roughly
1.5--2.0 tokens per character, and a naive character-truncation rule
silently exceeds the 8{,}192-token API limit on long Hebrew or
Chinese articles. For full-page Wikipedia content, inputs are
chunk-and-mean averaged: the article is split into 8{,}000-token
windows, each window embedded separately, and the resulting vectors
averaged. Lead-only inputs always fit in a single window and are
embedded directly.

For the cross-embedder robustness sweep
(\S\ref{sec:exp_015}), three independent multilingual embedders are
evaluated alongside the OpenAI baseline: Alibaba
\verb|text-embedding-v3| (1{,}024-d, Chinese-trained), Google
\verb|gemini-embedding-001| (3{,}072-d, US-trained on a different
corpus from OpenAI), and Yandex \verb|text-search-doc| (256-d,
Russian-trained). The three alt-embedders are accessed through their
respective REST APIs. Each embedder produces its own copy of
the response and anchor embeddings in per-embedder subtrees
(\verb|llm_embeddings_alt/|, \verb|processed_leads_alt/|) of the
event data tree, so cosine
matrices can be recomputed end-to-end in any embedder space.

\section{Cosine metric and ingroup pull}\label{sec:cosine_metric}

For a response embedding $\embed{r_\ell}$ to question $q$ in language
$\ell$, the diagonal ingroup pull is
\begin{equation}
  \ingrouppull(r, q, \ell)
    \;=\; \cossim\bigl(\embed{r_\ell},\, \embed{a_{q,\ell}}\bigr),
\end{equation}
where $a_{q,\ell}$ is the $\ell$-language Wikipedia lead paragraph
for question $q$. The full $3 \times 3$ response-language $\times$
anchor-language matrix supports a row-centred view that subtracts
the row mean and isolates ingroup pull from outgroup distance.
Throughout the thesis, ``raw cosine'' refers to the un-centred
cosine, ``row-centred'' to the within-row demeaned matrix, and
``ingroup pull'' specifically to the diagonal of the row-centred
matrix. Figure~\ref{fig:cosine_matrix_schematic} illustrates the
3$\times$3 layout for the RU--UK case.

\begin{figure}[!htbp]
  \centering
  \includegraphics[width=\linewidth]{cosine_matrix_schematic.png}
  \caption[Cosine matrix layout]{The 3$\times$3 cosine matrix for a
    single (provider, qid) cell, illustrated with the actual
    \texttt{claude-haiku-4-5} numbers on q02 (``what happened in
    Crimea in 2014?''). Diagonal cells (highlighted) are the
    matched-language ingroup pull; off-diagonal cells are the
    outgroup distance. The right panel summarises the row-centring
    transformation that isolates ingroup pull from outgroup distance.}
  \label{fig:cosine_matrix_schematic}
\end{figure}

\section{Statistical treatment}\label{sec:stats}

Uncertainty is reported at two levels. Within a (provider, question,
language) cell, the tables carry 95\% normal-approximation intervals
over the raw cosines; these bound sampling noise inside a cell only,
and because the ten samples within a question are draws from one
distribution, the effective sample size for generalising over
\emph{topics} is the number of questions, not the number of cosines.
Matrix-level claims therefore use a two-level cluster bootstrap:
providers and questions are resampled with replacement (2{,}000
replicates) and percentile intervals are reported for the
per-language means and their differences. Because four of the nine
questions share the annexation anchor, a stricter variant resamples
the six \emph{anchors} as blocks, carrying their attached
questions; both are reported in \S\ref{sec:exp_001}. The providers are not
independent draws --- they share training corpora --- and positive
dependence of this kind shrinks the effective sample size, so these
intervals are read as descriptive summaries that may under-cover
rather than as strict coverage guarantees.

The significance of the ingroup-pull association itself is tested
with the permutation null that transfers from
\citet{caliskan2017semantics}: within every (provider, question)
cell the response-language labels of the cosine matrix are randomly
relabelled, and the observed grand mean row-centred diagonal is
compared against $10{,}000$ permuted replicates
(\S\ref{sec:exp_001}). Stance-gap uncertainty uses the same
bootstrap design over providers and responses
(\S\ref{sec:exp_021}). Cluster--design association is scored by
adjusted mutual information rather than raw modal purity, because
purity is not comparable across axes with different numbers of
classes (\S\ref{sec:exp_016}). Rankings reported without intervals
--- the between-family ladder tables --- are descriptive and are
flagged as such where they appear.

\section{Language-subspace debiasing}\label{sec:debiasing}

The cosine pull on the diagonal can in principle reflect either
\emph{content} (the response and anchor share narrative framing) or
\emph{language} (the response and anchor share script and lexicon).
To separate these, the analysis projects out the language axis,
following the orthogonal-complement procedure of
\citet{bolukbasi2016man} adapted to multiple languages.

For language set \(\mathcal{L}\), the per-language centroid
\(\bvec{c}_\ell\) is computed by averaging the control embeddings
in language \(\ell\). The language axis $\langaxis$ is then the
orthonormal basis of the span of the centred centroid set
\(\{\bvec{c}_\ell - \bar{\bvec{c}}\}\), with $\bar{\bvec{c}}$ the
global control centroid and dimensionality at most
\(|\mathcal{L}| - 1\) (so two dimensions for the three-language
RU--UK case). Crucially, $\langaxis$ is estimated from controls
\emph{only}: the contested corpus is never used to fit the
projection, eliminating the circularity worry that one would
otherwise project away the very signal one is trying to measure.
That choice has a cost as well as a benefit, and both are stated:
the controls are Wikipedia article text while the projected objects
are chat completions, so the fitted axis is a language
$\times$ Wikipedia-register direction and may under-remove the
language component of \emph{responses}; and the axis removes at
most $|\mathcal{L}| - 1$ dimensions (two, for three languages) of a
1{,}536-dimensional space by construction, not by a variance
criterion, so ``survives the projection'' is a necessary rather
than sufficient test of content-drivenness. \S\ref{sec:exp_001_debiased}
brackets the estimate by refitting the axis from the response set's
own per-language centroids --- legitimate there because topic
composition is identical across languages by design.

Both response and anchor embeddings are projected onto the orthogonal
complement of $\langaxis$,
\begin{equation}
  \bvec{x} \;\mapsto\; \bvec{x} - \langaxis^\top \langaxis\, \bvec{x},
\end{equation}
before cosine recomputation
(Figure~\ref{fig:debiasing_geometry} illustrates the geometry). The
interpretation of the resulting
``debiased ingroup pull'' is the cosine similarity that survives
after the dimensions encoding pure language identity have been
removed: large positive values indicate content-driven cosine pull,
near-zero or negative values indicate that the original on-diagonal
pull was a lexical or scriptal artefact. The substantive cross-event
pattern of which language pairs survive debiasing and which collapse
is reported in \S\ref{sec:exp_006_debiased}.

\begin{figure}[!htbp]
  \centering
  \includegraphics[width=\linewidth]{debiasing_geometry.png}
  \caption[Language-axis debiasing geometry]{Geometric illustration
    of language-axis projection (schematic; the actual procedure
    runs in 1{,}536 dimensions). The grey plane is the language
    subspace fitted from the control centroids. A response embedding
    (filled point) decomposes into its projection onto that subspace
    --- the language component, drawn as the open circle --- and the
    orthogonal residual (dashed segment), which is the debiased
    coordinate the projection retains. Diagonal pull that survives
    the projection is content-driven; pull that collapses was
    lexical.}
  \label{fig:debiasing_geometry}
\end{figure}

\section{Topic-vs-language disentanglement}\label{sec:topic_vs_lang}

A second method for separating content-driven from language-driven
cosine pull bypasses the language-axis projection entirely.
For each (event, language) cell, the analysis computes two means:
\begin{enumerate}
  \item the \emph{on-topic} mean cosine, \(\overline{c}_{\mathrm{on}} =
    \mathbb{E}[\cossim(\bvec{r}, \bvec{a}_{q,\ell})]\), between
    responses and the event-specific Wikipedia lead anchor;
  \item the \emph{off-topic same-language baseline}
    \(\overline{c}_{\mathrm{off}} = \mathbb{E}[\cossim(\bvec{r},
    \bvec{u}_{\ell})]\), where $\bvec{u}_\ell$ is sampled
    \(N=200\) times from the universal-control corpus restricted to
    language \(\ell\).
\end{enumerate}
The \emph{topic lift}
\(\Delta_{\mathrm{topic}} = \overline{c}_{\mathrm{on}} -
\overline{c}_{\mathrm{off}}\)
is the topic-specific cosine pull above the language-similarity
floor. A topic lift of zero indicates that the response clusters
with the event anchor only because both are in the same language; a
positive topic lift indicates a genuine content-driven pull above
that floor. This metric does not depend on a fitted projection and
is therefore robust to assumptions about the dimensionality of the
language axis.

Figure~\ref{fig:topic_vs_language_method} shows the measured
(event, language) cells in this space; the collapse percentages that
pair with it are reported in \S\ref{sec:exp_006_debiased}.

\begin{figure}[!htbp]
  \centering
  \includegraphics[width=0.92\linewidth]{topic_vs_language_scatter.png}
  \caption[Topic-vs-language disentanglement]{The
    topic-vs-language disentanglement across all thirteen (event,
    language) cells of the five-conflict corpus, coloured by
    language. The $x$-axis is the off-topic
    same-language baseline, the $y$-axis is the on-topic cosine;
    vertical distance above the $y{=}x$ line is topic lift. Every
    cell sits well above the line; Hebrew combines the highest
    baseline with mid-range lift, the English cells combine low
    baselines with the highest lift.}
  \label{fig:topic_vs_language_method}
\end{figure}

\section{Paired-edit stance-axis projection}\label{sec:stance_axis}

Cosine to an own-language anchor answers \emph{how topically
similar is the response to its own-language encyclopedic
treatment?} It does not answer \emph{which side's framing
vocabulary does the response use?} The stance-axis projection is
built for the second question, adapting the bias-direction
construction of \citet{bolukbasi2016man}.

{\sloppy
For each conflict a \emph{paired-edit lexicon} is constructed:
twelve minimal-edit sentence pairs $(s^A_i, s^B_i)$ in English,
agent-flipped and structure-preserved, so that the members of a pair
share vocabulary and syntax and differ only in which party the
sentence assigns the action to. The stance axis is the unit vector
along the centroid difference of the two pole sets,
\begin{equation}
  \hat{\bvec{u}}_C \;=\;
  \frac{\bar{\bvec{s}}^A - \bar{\bvec{s}}^B}
       {\|\bar{\bvec{s}}^A - \bar{\bvec{s}}^B\|},
\end{equation}
and the signed stance projection of a response is the cosine of its
embedding onto that axis. For each (provider, language) cell the
per-response projections are averaged to a cell mean; the
cross-language gap for a provider is the maximum cell mean minus the
minimum cell mean across the conflict's languages; the per-conflict
headline is the mean of this gap across the provider matrix. Because
each lexicon fixes its own pole assignment, only the cross-language
\emph{gap} is comparable between families; the absolute level is
not.\par}

The five lexicons for the cross-conflict sweep (sixty pairs, $120$
seed sentences) live at
\verb|Code/src/babelbias/stance_lexicons.py| and were frozen in git
on 2026-05-18, before any response embedding was projected ---
building the vocabulary first prevents tuning the axis to a desired
outcome. Each axis is well-defined before any response touches it;
the seed separations range from $0.43$ (Israel--Palestine) to
$0.62$ (Falklands). The additional ladder-family and settled-pair
lexicons used in \S\ref{sec:exp_007} were committed later
(2026-06-26), together with the projection code and before any
ladder projection output existed, but without the multi-week
separation of the original freeze; the anti-tuning guarantee is
therefore strongest for the five-conflict sweep. Sensitivity to the particular pairs chosen is
checked with a leave-one-pair-out ablation reported alongside the
results (\S\ref{sec:exp_021}).

Two properties of the construction bound what the metric can claim.
First, the seed sentences are written in English, so projecting
non-English response embeddings onto the axis relies on the
embedder mapping translation-equivalent framing onto the same
direction. If the quality of that cross-lingual alignment varies
across language pairs, part of any between-family gap difference
could reflect alignment quality rather than framing behaviour; the
within-family cross-language gap, computed under a single axis, is
the safer comparison, and the between-family ranking should be read
with this in mind. Second, the cross-language gap is a
max-minus-min range statistic, and the expected range grows with
the number of cells it is taken over. Families measured in three
languages (the English control plus two contested languages) are
therefore favoured over two-language families (Taiwan strait,
Falklands) in the between-family ranking; comparisons between
families with equal language counts are unaffected.

\section{Anchor-free metrics for imaginary content}\label{sec:anchor_free}

The cosine-to-anchor metric requires a Wikipedia article on the
contested event. For the imaginary-conflict pilot
(\S\ref{sec:exp_004}) the contested event is fictional and no anchor
exists, so two anchor-free metrics are used.

\paragraph{Metric 1: stance coding.} An LLM-as-judge protocol scores
each response on a discrete \(-2/+2\) stance rubric: \(+2\) is a
response that clearly sides with the \emph{Russian} framing of the
fictional event, \(-2\) one that clearly sides with the
\emph{Ukrainian} framing, and \(0\) is neutral or hedged. Positive
therefore means Russian-ward; this is the reporting convention for
every stance-judge number in the thesis, with one flagged exception
(the ladder judge arm of \S\ref{sec:exp_007}, whose per-family axis
is signed toward the second-named party). The judge model is
Anthropic \verb|claude-sonnet-4-6|, deliberately disjoint from the
swept-model set so that no judged model is also the judge, and the
scoring prompt was fixed before any responses were seen.

\paragraph{Metric 2: within-model cross-lingual divergence.}
For each (provider, qid) pair, the metric averages
\(1 - \cossim(\bvec{r}_\ell, \bvec{r}_{\ell'})\) across all
language-pair samples within the same model. This is the within-model
analogue of the cosine ingroup pull and requires no anchor: it
measures only how differently the same model answers the same
question when prompted in different languages.

The two metrics are orthogonal: Metric~1 measures direction (which
side does the response lean toward?), Metric~2 measures magnitude
(how much do the answers differ across languages, regardless of
direction?). The discrepancy between them on the imaginary-conflict
pilot is itself a finding (\S\ref{sec:two_shapes}).

\section{Unsupervised clustering: EVoC}\label{sec:evoc}

A third axis-free check on the experimental design uses the EVoC
multi-resolution clustering algorithm \citep{mcinnes2024evoc} on the
full response embedding set, without supplying provider, question, or
language labels. The cluster purity along each design axis (qid,
language, model) at the algorithm's chosen resolution is reported as
a methodology defence: if the (qid \(\times\) language) cell
structure that the cosine analysis assumes \emph{is} the dominant
structure of the embedding space, EVoC should recover it
unsupervised.

In practice EVoC produces a cluster set whose modal purity is high
on both qid and language and low on provider (above chance by the
chance-adjusted AMI of \S\ref{sec:stats}, though only marginally
above the expected modal share under random labels --- see
Appendix~\ref{sec:appendix_evoc_purity}), across every
embedder space tested (\S\ref{sec:exp_016}). Anchor projection
(\S\ref{sec:exp_016_anchor_projection})~--- placing the Wikipedia
anchors as points inside the unsupervised EVoC fit~--- closes the
loop with the cosine analysis: for the questions whose anchors land
in a real cluster, anchors fall into the cluster of
their matched-language responses, which is the same statement as
positive ingroup pull, expressed structurally rather than
metrically.

\section{Tokenizer-overlap control}\label{sec:tokenizer_control}

The cross-language ingroup-pull pattern could in principle be a
tokenizer artefact. \citet{qi2023crosslingual} show that subword-vocabulary
overlap is the strongest predictor of cross-lingual factual
consistency in multilingual LMs, with BLOOM's RU--UK corpus-overlap
score reaching 0.76. If the embedder's tokenizer carves Russian and
Ukrainian into mostly-shared subwords, the embeddings cannot tell
them apart at the input layer, and any downstream cluster fusion is
a vocabulary artefact rather than a narrative finding.

The analysis pipeline therefore includes an explicit tokenizer-overlap
control (\S\ref{sec:exp_017}). For each of the three embedder
tokenizers that are publicly inspectable --- OpenAI
\verb|cl100k_base|, a Qwen tokenizer standing proxy for the Alibaba
embedder, and Cohere's multilingual tokenizer; the Google and Yandex
embedder tokenizers are not published --- pairwise corpus-overlap
scores between every pair of target languages are computed on the
778 control triples with full EN/RU/UK text available. The Wikipedia-anchor projection from
\S\ref{sec:evoc} is then re-run under the lowest-overlap embedder
for which response embeddings were collected --- Alibaba
\verb|text-embedding-v3| (Qwen-proxy overlap $0.42$); the Cohere
tokenizer scores lower still ($0.31$) but its embedding run was not
collected, so it serves as a tokenizer-only reference point. Fusion
is compared as the number of questions whose Ukrainian responses
join a Russian-anchored cluster at matched cluster count.
Convergent or divergent results between the two embedders constrain
the tokenizer-as-explanation reading.

\section{Cross-event scale-up}\label{sec:cross_event_scaleup}

The pipeline above is applied with no code changes to four further
contested-event families: Israel--Palestine (EN/HE/AR),
India--Pakistan (EN/HI/UR), Taiwan strait (EN/ZH), and the Falklands
War (EN/ES). The Falklands family is included deliberately as an
Oeberst-predicted near-null falsification anchor: \citet{oeberst2020wikipedia}
identify it as one of the conflicts where the ingroup-bias prediction
does not robustly hold in their human-rater data, so observing a
similarly weak signal in the LLM responses is a positive falsification
result rather than a project failure.

Each new event ships with a nine-question prompt bank mirroring the
\verb|ru_uk_core| structure (loaded terminology, event framing,
revolution / cause, legitimacy, attribution, sovereignty, two
forced-framing prompts, historical figure), Wikipedia anchor titles
in each target language resolved through the langlinks API, and the
universal-control language subset relevant for that event's
languages. Per-language prompts are translated using DeepL with a
secondary cross-check against an LLM-drafted version, and all
divergences are flagged for native-speaker review. The cross-event
sweep keeps temperature, token cap, and provider matrix identical to
the RU--UK headline; \(N\) was reduced to 5 samples per cell so the
total spend across all four new events fits comfortably under \$10
in measured LLM API costs (Appendix~\ref{sec:appendix_costs}).

The cross-event extension is the empirical surface for the
\emph{generalisation} question (RQ2 in \S\ref{sec:rqs}). Whether
the cosine ingroup pull, the language-axis debiasing pattern, the
topic-vs-language separation, and the EVoC structural finding all
generalise across event families is investigated in
\S\ref{sec:exp_006}.

\section{Event ladders and the four-arm measurement
  design}\label{sec:event_ladders}

The cross-event scale-up above varies the conflict while holding the
event date roughly constant. To vary time instead, the design is
extended with \emph{event ladders}: for a given bilateral dispute, a
set of focal events spanning as much of its documented history as
both parties' Wikipedia editions support. Within a ladder the
conflict family, the prompt bank structure, the provider set, and
the sampling protocol are all held fixed, so the event date is the
only systematically varying quantity. This is the controlled form of
the recency comparison that \citet{oeberst2020wikipedia}'s second
structural moderator implies, and it is the design behind
\S\ref{sec:exp_007}.

Eight ladders are built. Six cover contested disputes
(Russia--Ukraine, Israel--Palestine, India--Pakistan, China--Japan,
Greece--Turkey, Poland--Russia); two are single-event
\emph{settled baselines}, chosen as bilateral separations that
concluded without lasting contestation --- the 1993 Velvet Divorce
and the 1905 dissolution of the Norway--Sweden union. The settled
baselines are the design's floor. Each pairs two mutually
intelligible-adjacent languages with a genuine historical separation
but no live dispute, so whatever the metric reads on them is the
value that bilingual difference produces in the absence of
contestation. A metric that cannot separate contested from settled
pairs is measuring language, not framing.

Ladder anchors are resolved through the same \verb|langlinks| route
as the case-study anchors (\S\ref{sec:wiki_corpus}). Ladder prompts
mirror the \verb|ru_uk_core| bank structure and are
machine-translated into each family's languages. Apart from English,
and apart from the Russian and Ukrainian prompts inherited from the
case-study bank, no ladder prompt has been reviewed by a native
speaker; the corpus spans sixteen languages, several of which the
project has no internal capacity to read. The ladder arm is
therefore exploratory relative to the case-study sweep, and the
consequences of that are set out in \S\ref{sec:threat_mt}.

\paragraph{Four measurement arms.} The ladder corpus is scored under
four instruments rather than one, chosen so that the failure modes
of each are covered by at least one other.

\begin{enumerate}
  \item \emph{Anchor cosine} (\S\ref{sec:cosine_metric}) --- the
    row-centred diagonal pull, unchanged from the case-study
    protocol. Measures topical resemblance to the own-language
    Wikipedia treatment. Requires a published anchor in every
    language of the family.
  \item \emph{Paired-edit stance projection}
    (\S\ref{sec:stance_axis}) --- a per-family lexicon and axis,
    constructed as in the five-conflict sweep. Measures which side's
    framing vocabulary a response uses, controlled for topic
    overlap. As set out in \S\ref{sec:stance_axis}, only the
    cross-language \emph{gap} is comparable between families, and
    the between-family ranking inherits the English-seed-axis and
    range-statistic caveats stated there.
  \item \emph{Refusal rate} --- the proportion of completions in a
    cell that decline to answer, scored with the classifier from the
    Yandex sweep (\S\ref{sec:exp_011}). This arm exists to close off
    a specific alternative explanation: a diagonal that falls
    because models stop answering is a different phenomenon from a
    diagonal that falls because models answer differently.
  \item \emph{LLM-as-judge stance} (\S\ref{sec:anchor_free}) --- the
    stance rubric applied by a single judge model to two judged
    providers' responses on the most recent event in each family.
    Measures directional commitment
    independently of the embedding space, and so provides the only
    arm that shares no machinery with the other three --- though it
    shares its judge with the \S\ref{sec:anchor_free} stance code
    (\S\ref{sec:threat_single_judge}).
\end{enumerate}

The four arms are deliberately redundant on direction and
complementary on failure mode. Cosine and stance both run through
the embedding space and would move together under an embedder
artefact; the judge arm does not, and so is the check on that.
Cosine, stance, and the judge all score \emph{answered} responses;
the refusal arm is the check on whether the answered set is
representative. No arm is treated as decisive alone, and where they
disagree the disagreement is reported rather than resolved by
selection (\S\ref{sec:four_instruments}).
